\answer{ex:04:1}{(a)~$\approx1.1\cdot10^{11}$ bits/J. (b)~$\log_2\sqrt{10}\approx1.7$ bit contra $81$: $\approx50$ vezes.}
\answer{ex:04:2}{$r^*=(P_0/E_{\text{disp}})\ln\bigl(1/(r^*\Delta t)\bigr)$, $P_0/E_{\text{disp}}=1.16\,\text{s}^{-1}$: $r^*\approx6$ disparos por segundo, $7.4\cdot10^{10}$ bits/J.}
\answer{ex:04:3}{$\mathcal I=\frac{N}{1+(N-1)c}=9.9$ (limite $10$) contra $\mathcal I=\frac{N}{1-c}=1111$: $\approx110$ vezes; a correlação só é nociva quando se parece com o sinal.}
\answer{ex:04:4}{$\theta\approx68.7$ e $19.9$ (em unidades de $w$), $m=19$ e $m=10$: o destinatário rápido precisa de quase metade das entradas coincidentes, embora sua entrada média seja cinco vezes menor.}
\answer{ex:04:5}{$U\tau_{\text{rec}}\ge0.725\,\text{s}$ e $\tau_{\text{rec}}(1-2U)\le0.05\,\text{s}$, o que exige $U\ge0.483$; o menor $\tau_{\text{rec}}=0.725\,\text{s}$ com $U=1$ ($1.45\,\text{s}$ com $U=0.5$).}
\answer{ex:04:6}{(a)~$\theta\,e/(e-1)\approx1.58\,\theta$ para quaisquer $\tau$ e $\theta$. (b)~$14$ disparos.}
\answer{ex:04:7}{(a)~$1.62$ bit por palavra: $0.77$ do número de disparos, $0.84$ de sua disposição. (b)~$1.13$ e $1.59$ bit: com $30$ palavras, a estimativa perde quase um terço.}
